% ---------------------------------------------------------------
%  Section X -- Conclusion
% ---------------------------------------------------------------

We presented IRIDIUM, a provenance-aware urban mobility platform for
Azerbaijani cities that unifies a real-data dynamic digital twin
$G=(V,E,W)$, a modular analytics suite, and a federated learning pathway
under a strict no-synthetic-substitution policy.
Every API response carries a machine-readable \texttt{data\_status} field
and a provenance list $\Pi$ that completely characterizes the data lineage
of the result.
The mathematical formulations for multi-objective routing
(Eq.~\ref{eq:routing}), ST-GNN forecasting
(Eqs.~\ref{eq:forecast}--\ref{eq:gcn}), Mobility Equity Score
(Eqs.~\ref{eq:mes}--\ref{eq:mes_var}), anomaly severity
(Eqs.~\ref{eq:anomaly_severity}--\ref{eq:fpr}),
FedAvg aggregation (Eq.~\ref{eq:fedavg}), and federated convergence
(Eq.~\ref{eq:fl_convergence}) are fully specified and consistent with
the implemented codebase, enabling future work to instantiate each
component without ambiguity.

Baseline evaluation confirmed API semantic correctness, graceful
degradation under unconfigured sources, end-to-end latency of
2--55\,ms (median) under light load, and successful real weather
integration for Baku and Quba.
The empirically measured temperature differential of
$\overline{\Delta T} = 2.50 \pm 0.82\,\text{\textdegree{}C}$ is
consistent with the expected superposition of orographic lapse
($\approx 3.8\,\text{\textdegree{}C}$ for the \SI{582}{m} elevation
difference) and Caspian Sea thermal buffering, confirming that multi-city
atmospheric data are both retrievable and physically coherent.
Three precipitation events were identified in the 72-hour Baku series,
yielding approximately 16\,mm cumulative precipitation and confirming
the operational readiness of the weather input for anomaly severity scoring.

Realistic next steps, in priority order, are as follows.

\begin{enumerate}
  \item Validating the OSM network import on the full Azerbaijan PBF
        extract, verifying routing on representative Baku
        origin-destination pairs, and benchmarking Dijkstra-based routing
        latency against production load targets.
  \item Negotiating GTFS static and GTFS-RT data-sharing agreements with
        Baku Metro and BakuBus to enable multimodal routing and
        transit-aware congestion forecasting.
  \item Integrating a licensed or open traffic-speed feed to activate
        dynamic edge state $\mathbf{d}_e(t)$ and evaluating the heuristic
        forecaster and the ST-GNN (Table~\ref{tab:hyperparams}) against
        held-out observations, reporting MAE, RMSE, MAPE, and sMAPE
        (Eqs.~\ref{eq:mae}--\ref{eq:mape}).
  \item Assembling a district-level accessibility indicator dataset for
        Baku and evaluating $\MES_d$ and $\Delta_{\MES}$ against external
        benchmarks, with sensitivity analysis over indicator weights $w_m$.
  \item Deploying the federated learning pathway with local training nodes
        at partner operator sites and characterizing convergence
        (Eq.~\ref{eq:fl_convergence}) and the privacy-utility trade-off
        (Eq.~\ref{eq:dp_noise}) empirically.
\end{enumerate}

The repository, mathematical specification, provider inventory, and
data-status protocol presented in this paper constitute a transparent,
auditable foundation for these extensions, without overclaiming the
current maturity of the system.
The platform design is intended to serve both as a practical baseline
for Azerbaijani mobility deployments and as a reproducible reference
architecture for urban mobility systems operating in data-scarce,
governance-constrained environments.
